\documentclass[border=2pt]{standalone}
% Estilo comum dos fluxogramas do Caderno C3 (cores do modelo do MOP)
\usepackage[T1]{fontenc}
\usepackage[utf8]{inputenc}
\usepackage{helvet}
\renewcommand{\familydefault}{\sfdefault}
\usepackage{amssymb}
\usepackage{tikz}
\usetikzlibrary{positioning, arrows.meta, shapes.geometric, shapes.symbols, shapes.misc,
  calc, fit, backgrounds, decorations.pathmorphing, patterns, shadows}
% Paleta do modelo do caderno
\definecolor{cTeal}{HTML}{00747A}
\definecolor{cNavy}{HTML}{0B3A5A}
\definecolor{cBlue}{HTML}{0A73BA}
\definecolor{cBlueBg}{HTML}{EBF4F9}
\definecolor{cTealBg}{HTML}{E3F1F2}
\definecolor{cTealMid}{HTML}{9CCFD2}
\definecolor{cGrayBg}{HTML}{F2F2F2}
\definecolor{cGrayLine}{HTML}{BFBFBF}
\definecolor{cGrayText}{HTML}{595959}
\definecolor{cText}{HTML}{1A1A1A}
\definecolor{cAmber}{HTML}{B7791F}
\definecolor{cAmberBg}{HTML}{FDF3DC}
\definecolor{cGreen}{HTML}{3D7D3A}
\definecolor{cGreenBg}{HTML}{E8F2E4}
\color{cText}
\tikzset{
  every node/.style={font=\footnotesize, text=cText},
  etapa/.style={draw=cTeal, fill=cTealBg, rounded corners=2pt, align=center,
     inner sep=3.5pt, minimum height=0.85cm, text width=#1, line width=0.6pt},
  etapa/.default={3.2cm},
  marco/.style={draw=cNavy, fill=cNavy, rounded corners=2pt, align=center,
     inner sep=3.5pt, minimum height=0.85cm, text width=#1,
     font=\footnotesize\bfseries, text=white},
  marco/.default={3.2cm},
  marcoteal/.style={draw=cTeal, fill=cTeal, rounded corners=2pt, align=center,
     inner sep=3.5pt, minimum height=0.85cm, text width=#1,
     font=\footnotesize\bfseries, text=white},
  marcoteal/.default={3.2cm},
  externo/.style={draw=cGrayLine, fill=cGrayBg, rounded corners=2pt, align=center,
     inner sep=3.5pt, minimum height=0.85cm, text width=#1},
  externo/.default={3.2cm},
  produto/.style={draw=cTeal!70, fill=white, rounded corners=1pt, align=center,
     inner sep=2.5pt, minimum height=0.6cm, text width=#1, font=\scriptsize},
  produto/.default={2.4cm},
  decisao/.style={draw=cAmber, fill=cAmberBg, diamond, aspect=2, align=center,
     inner sep=1pt, font=\scriptsize, line width=0.6pt},
  correcao/.style={draw=cAmber, fill=cAmberBg, rounded corners=2pt, align=center,
     inner sep=3pt, minimum height=0.8cm, text width=#1, font=\scriptsize},
  correcao/.default={2.6cm},
  transversal/.style={draw=cBlue, fill=cBlueBg, dashed, rounded corners=2pt, align=center,
     inner sep=3.5pt, text width=#1, line width=0.6pt},
  transversal/.default={6cm},
  seta/.style={-{Stealth[length=2.2mm,width=1.8mm]}, line width=0.8pt, draw=cNavy},
  setaleve/.style={-{Stealth[length=1.8mm,width=1.5mm]}, line width=0.6pt, draw=cGrayText},
  setatracejada/.style={-{Stealth[length=2mm]}, line width=0.7pt, dashed, draw=cGrayText},
  rotulo/.style={font=\scriptsize\itshape, text=cGrayText, fill=white, inner sep=1pt},
  fase/.style={signal, signal to=east, signal from=west, draw=cTeal, fill=cTealBg,
     align=center, font=\scriptsize\bfseries, minimum height=1.05cm, text width=#1,
     inner sep=2pt},
  fase/.default={1.9cm},
  raia/.style={draw=cGrayLine, fill=cGrayBg, minimum width=#1, minimum height=0.9cm,
     anchor=north west},
  tituloraia/.style={font=\scriptsize\bfseries, text=cNavy, align=center, text width=1.75cm},
  cab/.style={font=\scriptsize\bfseries, text=white, fill=cNavy, align=center,
     inner sep=3pt, text width=#1, minimum height=0.6cm},
  cab/.default={3cm},
  bloco/.style={draw=cGrayLine, fill=cGrayBg, rounded corners=3pt},
}
% sem hifenização nas caixas dos fluxogramas
\AtBeginDocument{\hyphenpenalty=10000 \exhyphenpenalty=10000 \sloppy}

\begin{document}

\begin{tikzpicture}[x=1cm,y=1cm,
   num/.style={circle, fill=petroleo, text=white, font=\fonte{8}{9}\bfseries, inner sep=0pt,
     minimum size=0.44cm},
   traco/.style={draw=petroleo, line width=0.8pt},
   cano/.style={draw=petroleo, line width=2pt},
   nota/.style={font=\fonte{7.5}{8.8}, text=cinza80}]
  % solo e aquífero
  \fill[cinza5] (0,0) rectangle (15.6,-3.5);
  \fill[petroleo12] (0,-2.35) rectangle (15.6,-3.5);
  \fill[pattern=dots, pattern color=petroleo!60] (0,-2.35) rectangle (15.6,-3.5);
  \draw[cinza60, line width=0.5pt] (0,-2.35) -- (15.6,-2.35);
  \node[nota, anchor=west, fill=petroleo12, inner sep=1.5pt] at (2.3,-2.95) {aquífero (água subterrânea)};
  \node[nota, anchor=west] at (2.3,-1.3) {solo e rocha};
  \draw[cinza60, line width=1.2pt] (0,0) -- (15.6,0);
  % poço
  \fill[white] (1.25,0.3) rectangle (1.75,-3.1);
  \draw[traco] (1.25,0.3) -- (1.25,-3.1) -- (1.75,-3.1) -- (1.75,0.3);
  \draw[traco, dash pattern=on 2pt off 1.5pt] (1.25,-2.45) -- (1.25,-3.05);
  \draw[traco, dash pattern=on 2pt off 1.5pt] (1.75,-2.45) -- (1.75,-3.05);
  \fill[petroleo12] (1.27,-1.9) rectangle (1.73,-3.08);
  \filldraw[traco, fill=white] (1.36,-2.9) rectangle (1.64,-2.45);
  \draw[traco] (1.5,-2.45) -- (1.5,0.3);
  \filldraw[traco, fill=white] (1.1,0.3) rectangle (1.9,0.55);
  \node[num] at (0.65,0.85) {1};
  % adutora até a casa química
  \draw[cano] (1.9,0.42) -- (3.7,0.42);
  % casa química
  \filldraw[traco, fill=white] (3.7,0) rectangle (5.6,1.35);
  \filldraw[traco, fill=white] (3.55,1.35) -- (4.65,1.95) -- (5.75,1.35) -- cycle;
  \node[nota, align=center] at (4.65,0.68) {cloração\\ e controle};
  \node[num] at (4.65,2.35) {2};
  % ligação ao reservatório
  \draw[cano] (5.6,0.42) -- (7.05,0.42) -- (7.05,2.75) -- (7.4,2.75);
  % reservatório elevado
  \draw[traco] (7.6,0) -- (7.75,2.2);
  \draw[traco] (9.0,0) -- (8.85,2.2);
  \draw[traco, line width=0.5pt] (7.68,1.1) -- (8.92,1.1);
  \fill[petroleo12] (7.4,2.2) rectangle (9.2,3.05);
  \draw[traco] (7.4,2.2) rectangle (9.2,3.6);
  \draw[traco, line width=0.5pt] (7.4,3.05) -- (9.2,3.05);
  \node[num] at (8.3,4.0) {3};
  % rede de distribuição
  \draw[cano, line width=1.6pt] (8.3,2.2) -- (8.3,-0.4) -- (14.9,-0.4);
  \node[num] at (8.75,-0.8) {4};
  % domicílios com hidrômetro
  \foreach \x in {10.9, 12.65, 14.4} {
    \draw[traco, line width=1.1pt] (\x-0.75,-0.4) -- (\x-0.75,0.15) -- (\x-0.45,0.15);
    \node[circle, draw=petroleo, fill=white, line width=0.7pt, inner sep=1.4pt] at (\x-0.75,0.15) {};
    \filldraw[traco, fill=white] (\x-0.45,0) rectangle (\x+0.45,0.8);
    \filldraw[traco, fill=white] (\x-0.6,0.8) -- (\x,1.25) -- (\x+0.6,0.8) -- cycle;
    \draw[traco, line width=0.6pt] (\x-0.15,0) rectangle (\x+0.15,0.45);
  }
  \node[num] at (10.15,0.75) {5};
  \node[num] at (15.25,0.55) {6};
  % legenda
  \begin{scope}[shift={(9.55,4.35)}]
    \foreach \n/\t [count=\k from 0] in {1/Poço tubular profundo com bomba submersa,
        2/{Casa química: cloração da água}, 3/Reservatório, 4/Rede de distribuição,
        5/Ligação domiciliar com hidrômetro, 6/Domicílio atendido} {
      \node[num, minimum size=0.36cm, font=\fonte{7}{8}\bfseries] at (0.18,{-0.2-\k*0.42}) {\n};
      \node[anchor=west, font=\fonte{8}{9.4}] at (0.42,{-0.2-\k*0.42}) {\t};
    }
  \end{scope}
  % anotações
  \node[nota, align=left, anchor=west] at (1.95,-0.6) {outorga de uso e análises\\ físico-químicas e bacteriológicas};
  \node[nota, align=center] at (12.55,-1.05) {operação e manutenção pela entidade gestora};
\end{tikzpicture}

\end{document}
